% SCOPE: 3.1 operational meta-block protocol + two-stage SFT priming; 3.2 the
% PMI-shift reward (two teacher-forced contexts, evidence margin, clipped shift,
% asymmetric save/derail bonus, region-routed token advantages, group centering);
% 3.3 the matched-base twin that isolates the meta package.
% Grounded in src/training/dcpo_pmi_shift.py and src/training/dcpo_region.py.
\section{Method}
\label{sec:method}

Our method has three parts. First, an operational protocol turns metacognition
into an explicit object that we can locate in the token stream. Second, a reward
scores each metacognitive block by the belief update it causally produces in the
model's own answer preferences. Third, a byte-identical control model isolates
the contribution of this reward from every other ingredient of the training
pipeline.

\subsection{An operational, localizable definition of metacognition}
\label{sec:method:protocol}

We make metacognition an explicit object in the model's output rather than a
trait we infer after the fact. At any point inside its reasoning, the model may
open a short delimited block, state three things, close the block, and resume
solving:
\begin{center}
\metaopen{} \texttt{confidence:} $\hat{p}$ \quad
\texttt{assessment:} \emph{(a judgment of the current attempt)} \quad
\texttt{action:} \emph{(verify / redirect / continue)} \metaclose{}
\end{center}
The opening marker \metaopen{} and the closing marker \metaclose{} are dedicated
tokens added to the vocabulary. Inside the block the model writes a confidence
number, an assessment of how the partial solution is going, and a committed next
action such as re-checking the current answer or switching to a different
approach. We call this representation executable because it exposes quantities we
can read directly. The block sits at a known position in the trace. Its content
is machine-parsable. The token context just before the block and the token
context just after it are both well defined, so we can measure how the block
changes what the model believes downstream.

Two supervised stages install this protocol before any reinforcement learning.
The first stage teaches form and placement. It trains on reasoning traces that
contain well-formed blocks embedded inside the chain of thought, so the model
reliably emits syntactically valid, correctly delimited blocks in the middle of
a solution. The second stage teaches function. It adds traces in which the block
is followed by a real behavioral consequence, such as abandoning a failing
approach or re-deriving a candidate answer, so that a block precedes an actual
change in the solving trajectory rather than sitting inert. We prime both form
and function first because reinforcement learning amplifies behaviors the policy
already produces rather than installing new ones \citep{gandhi2025cognitive}.
After priming, the reinforcement stage decides which of these behaviors to
strengthen and where they actually help.

\subsection{PMI-shift: rewarding a block by the evidence it moves}
\label{sec:method:pmishift}

A metacognition block should earn reward for the belief update it causes. It
should not be rewarded merely for existing, and not for the correctness of the
rollout it happens to sit in. This is a form of self-distillation with an unusual
target: the model already carries a latent preference between the correct answer
and a plausible wrong one in its own likelihoods, and we reward it for writing
explicit text that surfaces that preference and moves it in the right direction.
The equations below make this principle precise. We name the resulting reward
\pmishift{}, after the pointwise-mutual-information contrast between two candidate
answers that it is built on.

To measure the update, we read the model's answer preference at two contexts of
the same rollout and take the difference. Let $x$ be the problem and let
$y_{<b}$ be the reasoning the model produced strictly before it opened the block.
The context just before the block is the prompt plus that prefix, and the context
just after the block adds the block $b$ itself:
\begin{equation}
\mathrm{ctx}_{\mathrm{open}} = (x,\, y_{<b}),
\qquad
\mathrm{ctx}_{\mathrm{close}} = (x,\, y_{<b},\, b).
\label{eq:contexts}
\end{equation}
We compare two candidate final answers at each context. The first is the gold
answer $a^{\star}$, rendered as a boxed answer string. The second is a decoy
$\tilde{a}$: a plausible wrong answer produced by perturbing the gold string with
rule-based edits such as numeric offsets, sign flips, and numerator--denominator
swaps. The decoy is generated deterministically and checked to be non-equivalent
to the gold under a symbolic answer-equivalence checker, falling back to string
and float inequality when the checker itself errors. A frozen reference model
$\pi_{\mathrm{ref}}$, a copy of the policy that is never updated, reads the
probability of each candidate answer at each context. Reading here means
teacher-forcing: we feed the reference model a fixed answer string and record the
log-probability it assigns, with no sampling. Let
$\log p_{\mathrm{ref}}(a \mid \mathrm{ctx})$ be that teacher-forced log-probability
of answer string $a$ at context $\mathrm{ctx}$, summed over the answer tokens
where the gold and decoy tokenizations differ. The evidence margin at a context
is the reference model's log-odds of gold over decoy there:
\begin{equation}
m(\mathrm{ctx})
= \log p_{\mathrm{ref}}(a^{\star} \mid \mathrm{ctx})
- \log p_{\mathrm{ref}}(\tilde{a} \mid \mathrm{ctx}).
\label{eq:margin}
\end{equation}
Reading the margin at both contexts gives $m(\mathrm{ctx}_{\mathrm{open}})$, the
gold-versus-decoy evidence before the model reads its own block, and
$m(\mathrm{ctx}_{\mathrm{close}})$, the evidence after. These are four
teacher-forced readings in all: gold and decoy, each at the open and close
context. Their difference is the block's marginal contribution to belief:
\begin{equation}
\mathrm{SHIFT}
= m(\mathrm{ctx}_{\mathrm{close}}) - m(\mathrm{ctx}_{\mathrm{open}}).
\label{eq:shift}
\end{equation}
Everything else about the rollout stays fixed across the two contexts: the
prompt, the prefix, both candidate answers, and the reference model. So
$\mathrm{SHIFT}$ isolates the effect of the block itself.

The reward turns the shift into a bounded score and adds an asymmetric bonus for
genuine reversals of preference:
\begin{equation}
R_{\mathrm{shift}}
= s \cdot \mathrm{clip}(\mathrm{SHIFT},\, -c,\, c)
+
\begin{cases}
+\beta_{\mathrm{save}}
  & \text{if } m(\mathrm{ctx}_{\mathrm{open}}) < 0
    \text{ and } m(\mathrm{ctx}_{\mathrm{close}}) > 0
    \quad(\textsc{save}),\\[2pt]
-\beta_{\mathrm{derail}}
  & \text{if } m(\mathrm{ctx}_{\mathrm{open}}) > 0
    \text{ and } m(\mathrm{ctx}_{\mathrm{close}}) < 0
    \quad(\textsc{derail}),\\[2pt]
0 & \text{otherwise,}
\end{cases}
\label{eq:reward}
\end{equation}
with scale $s = 1$, clip bound $c = 2$, $\beta_{\mathrm{save}} = 1.0$, and
$\beta_{\mathrm{derail}} = 2.0$. A \textsc{save} is the event of interest: the
block flipped a decoy-leaning preference to a gold-leaning one. A \textsc{derail}
is the opposite failure, where the block talked the model out of a correct
leaning. The derail penalty is strictly larger than the save bonus, so harmful
metacognition costs more than useful metacognition earns. The clip keeps a single
large log-odds swing from dominating the objective. Degenerate cases fail closed
to $R_{\mathrm{shift}} = 0$ in one rule: a non-finite reading, an empty block
whose two contexts coincide, a decoy that collapses onto the gold string, or a
gold and decoy whose tokenizations cannot be length-matched all yield no shift
credit rather than a fabricated one.

The scalar reward becomes token-level advantages routed by region, so that
answer quality and block quality cannot game each other. Correctness reward
$R_{\mathrm{corr}}\in\{-1,+1\}$ flows
only onto the final-answer tokens, and $R_{\mathrm{shift}}$ flows only onto the
tokens inside the meta block. Because the two signals never share a token,
improving the answer cannot inflate the block reward and vice versa. These are
the two scientifically load-bearing heads: a correctness head on the answer
tokens and the shift head on the meta-content tokens. For each head, the scalar
is centered by subtracting the mean reward over the group of rollouts that share
a prompt, with no division by the group's variance, and the centered value is
broadcast uniformly onto every token of its own region:
\begin{equation}
A_{t}
= w_{\mathrm{corr}}\, \hat{A}_{\mathrm{corr}}\, \mathbb{1}[t \in \textsc{answer}]
+ w_{\mathrm{meta}}\, \hat{A}_{\mathrm{shift}}\, \mathbb{1}[t \in \textsc{meta}]
+ A^{\mathrm{aux}}_{t},
\qquad
\hat{A}^{(i)}_{\bullet} = R^{(i)}_{\bullet} - \frac{1}{|\mathcal{G}_{\bullet}|}\sum_{j\in\mathcal{G}_{\bullet}} R^{(j)}_{\bullet},
\quad \bullet\in\{\mathrm{corr},\mathrm{shift}\}.
\label{eq:routing}
\end{equation}
Here $\mathbb{1}[\cdot]$ is 1 when token $t$ lies in the named region and 0
otherwise, $\hat{A}^{(i)}_{\bullet}$ is the centered advantage for head $\bullet$,
and $\mathcal{G}_{\bullet}$ is the set of rollouts over which that head is
centered; for rollout $i$, $\hat{A}_{\mathrm{corr}}$ and $\hat{A}_{\mathrm{shift}}$
abbreviate $\hat{A}^{(i)}_{\mathrm{corr}}$ and $\hat{A}^{(i)}_{\mathrm{shift}}$.
Also $w_{\mathrm{corr}} = 1.0$, and $w_{\mathrm{meta}}$ is
warmed up linearly from $0$ to $0.8$ over the first $80$ steps. The shift head is
centered only over rollouts that actually emitted a scored block, so
non-emitting rollouts neither dilute the baseline nor receive a spurious
advantage. The auxiliary term $A^{\mathrm{aux}}_{t}$ collects a small set of
maintenance heads whose job is to keep the block channel clean, not to drive the
science, each routed onto its own tokens:
\begin{center}
\begin{tabular}{@{}lll@{}}
\toprule
Maintenance head & What it scores & Weight \\
\midrule
calibration & squared error of the stated confidence versus outcome & $0.30$ \\
format      & whether the block is well-formed and properly closed  & $0.35$ \\
emission    & whether the response opens a block at all             & $0.10$ \\
length cost & response length as a fraction of the token budget     & $0.08$ \\
truncation  & a block left open when the response is cut off        & $0.30$ \\
\bottomrule
\end{tabular}
\end{center}
The calibration head reads the confidence-number tokens, which are carved out of
the shift region so that writing a confidence the continuation then echoes cannot
earn shift credit. To keep weak and strong signals comparable, the auxiliary
heads are rescaled so their typical advantage magnitude tracks a running average
of the correctness head's, which prevents the small shift signal from being
buried under correctness.

\subsection{The matched-base twin that isolates the mechanism}
\label{sec:method:twin}

Any accuracy gap between our model and an off-the-shelf reinforcement-learning
baseline is hard to attribute, because such a baseline differs from our model in
many ways at once: the fine-tuning data, the prompt format, the token budget
spent on blocks, and the optimizer settings. We therefore train a matched-base
twin that differs from the treatment model in one thing only, the meta package
defined below.

The two models are held identical everywhere except the package. They train on
the same problems in the same order. The twin's fine-tuning data is
meta-stripped: every block is removed while the surrounding reasoning is kept
verbatim, so both models see the same mathematical content. Every non-meta
setting is identical, including the optimizer, the learning-rate schedule, the
group size, the sampling parameters, the sequence budget, and the number of
training steps. We enforce this with a byte-level audit of the two configurations
that whitelists only the meta-specific keys. What differs is the meta package as
a whole. The package is the priming protocol, the \pmishift{} reward, and the
maintenance heads together. The twin runs plain reinforcement learning with the
correctness reward alone: no meta tokens, no region routing, no shift head, and
none of the maintenance heads.

This contrast estimates the causal effect of the package. A between-model
comparison is necessary here. The alternative, comparing responses that emit a
block against responses that do not within a single model, would condition on a
decision the model itself makes, and rollouts that choose to emit a block may
differ systematically from those that do not. The matched-base contrast avoids
this by comparing two models on the same, full evaluation set. It estimates the
effect of the whole package rather than the shift reward alone. The experiments
further separate the package into the contribution of the priming stages and the
contribution of the reward, by comparing a meta-primed model trained with plain
reinforcement learning against the full treatment.
